\chapter{聚类分析}
前两章以变量为分析对象，本章起以观察对象（样品）为分析对象。聚类分析的目的是按相似程度将观察对象归为若干类。它与第4章判别分析的区别在于事先是否已知类别：聚类分析没有已知的类别可供对照，类别由数据得出，属于探索性分析；判别分析中训练样品的类别已知，目的是建立判别新样品的规则。

聚类分析的每一步都涉及选择：纳入哪些变量、是否标准化、采用何种距离、类间距离如何定义、分为几类。每一项选择都会影响结果，且没有统计检验能够判定哪种选择正确。因此，学习本章应掌握各项选择的依据，以及改变选择对结果可能产生的影响。

此外，聚类算法对任何数据都会给出分类结果，即使数据中并不存在群组结构。聚类结果需要结合专业知识解释，并宜通过稳定性分析或新资料加以验证。
\section{聚类分析}
聚类分析（cluster analysis）是一种探索性分析，将相似的观测或指标归类。分为Q型（样品）聚类和R型（指标）聚类两类。关键在于如何将相似性数量化，即相似系数。
\begin{center}\begin{tabular}{lll}\toprule
{\bfseries 测量尺度}&{\bfseries 英文}&{\bfseries 变量类型}\\\midrule
间隔尺度&Scale&数值变量\\
有序尺度&Ordinal&有序变量\\
名义尺度&Nominal&无序分类变量\\\bottomrule
\end{tabular}\end{center}
\Note{默认采用间隔尺度。}
\section{样品间相似性的度量}
每个样品可看成$p$维空间中的一个点，用距离度量样品间的相似程度。常用距离有欧氏距离、马氏距离、配合距离等。
\subsection{明氏距离}
\begin{align}
d_{ij}(q)&=\left[\sum_{k=1}^p|x_{ik}-x_{jk}|^q\right]^{1/q},&&\text{明氏距离},\\
d_{ij}(2)&=\left[\sum_{k=1}^p(x_{ik}-x_{jk})^2\right]^{1/2},&&\text{欧氏距离},\\
d_{ij}(1)&=\sum_{k=1}^p|x_{ik}-x_{jk}|,&&\text{绝对值距离},\\
d_{ij}(\infty)&=\max_k|x_{ik}-x_{jk}|,&&\text{切比雪夫距离}.
\end{align}
其中$q\geqslant1$。$q<1$时三角不等式不再成立，上式不是距离。
\Example{体重、胸围和往返跑数据}
\par\addvspace{5pt}\begin{center}\begin{minipage}{\linewidth}\centering
\small
\captionof{table}{体重、胸围与往返跑数据}
\begin{tabular}{lrrr}\toprule
{\bfseries 编号 }&{\bfseries  tz }&{\bfseries  xw }&{\bfseries  wfp}\\\midrule
1 & 72.0 & 87.0 & 14.5\\
2 & 63.4 & 81.0 & 11.4\\
3 & 54.1 & 82.5 & 12.2\\
4 & 79.9 & 95.0 & 10.9\\
5 & 74.8 & 88.9 & 13.9\\
\bottomrule\end{tabular}
\end{minipage}\end{center}

\begin{code}
da \ROperator{<-} \RFunction{data.frame}(tz \ROperator{=} \RFunction{c}(\RNumber{72.0, 63.4, 54.1, 79.9, 74.8}),
                 xw \ROperator{=} \RFunction{c}(\RNumber{87.0, 81.0, 82.5, 95.0, 88.9}),
                 wfp \ROperator{=} \RFunction{c}(\RNumber{14.5, 11.4, 12.2, 10.9, 13.9}))
x \ROperator{<-} \RFunction{cbind}(\RFunction{scale}(da\ROperator{\$}tz), \RFunction{scale}(da\ROperator{\$}xw), \RFunction{scale}(da\ROperator{\$}wfp))
x
\end{code}
\par\addvspace{5pt}\begin{center}\begin{minipage}{\linewidth}\centering
\small
\captionof{table}{标化后的数据}
\begin{tabular}{lrrr}\toprule
{\bfseries 编号 }&{\bfseries  tz }&{\bfseries  xw }&{\bfseries  wfp}\\\midrule
1 & 0.3103518 & 0.0215711 & 1.2273947\\
2 & -0.5342765 & -1.0569831 & -0.7543363\\
3 & -1.4476536 & -0.7873446 & -0.2429219\\
4 & 1.0862313 & 1.4596433 & -1.0739704\\
5 & 0.5853470 & 0.3631132 & 0.8438339\\
\bottomrule\end{tabular}
\end{minipage}\end{center}

\begin{code}
d \ROperator{<-} \RFunction{dist}(x, \RString{"euclidean"}); d
d \ROperator{<-} \RFunction{dist}(x, \RString{"manhattan"}); d
d \ROperator{<-} \RFunction{dist}(x, \RString{"minkowski"}, p \ROperator{=} \RNumber{1}); d
d \ROperator{<-} \RFunction{dist}(x, \RString{"maximum"}); d
\end{code}
\par\addvspace{5pt}\begin{center}\begin{minipage}{\linewidth}\centering
\small
\captionof{table}{欧氏距离}
\begin{tabular}{lrrrrr}\toprule
 &{\bfseries  1 }&{\bfseries  2 }&{\bfseries  3 }&{\bfseries  4 }&{\bfseries  5}\\\midrule
1 & 0.0000000 &  &  &  & \\
2 & 2.4091355 & 0.0000000 &  &  & \\
3 & 2.4303823 & 1.0809752 & 0.0000000 &  & \\
4 & 2.8224673 & 3.0102525 & 3.4871433 & 0.0000000 & \\
5 & 0.5825739 & 2.4133749 & 2.5763701 & 2.2652232 & 0.0000000\\
\bottomrule\end{tabular}
\end{minipage}\end{center}
\par\addvspace{5pt}\begin{center}\begin{minipage}{\linewidth}\centering
\small
\captionof{table}{绝对值距离（同 $q=1$ 的明氏距离）}
\begin{tabular}{lrrrrr}\toprule
 &{\bfseries  1 }&{\bfseries  2 }&{\bfseries  3 }&{\bfseries  4 }&{\bfseries  5}\\\midrule
1 & 0.0000000 &  &  &  & \\
2 & 3.9049135 & 0.0000000 &  &  & \\
3 & 4.0372376 & 1.6944301 & 0.0000000 &  & \\
4 & 4.5153168 & 4.4567682 & 5.6119212 & 0.0000000 & \\
5 & 1.0000983 & 4.1378901 & 4.2702142 & 3.5152185 & 0.0000000\\
\bottomrule\end{tabular}
\end{minipage}\end{center}
\par\addvspace{5pt}\begin{center}\begin{minipage}{\linewidth}\centering
\small
\captionof{table}{切比雪夫距离}
\begin{tabular}{lrrrrr}\toprule
 &{\bfseries  1 }&{\bfseries  2 }&{\bfseries  3 }&{\bfseries  4 }&{\bfseries  5}\\\midrule
1 & 0.0000000 &  &  &  & \\
2 & 1.9817310 & 0.0000000 &  &  & \\
3 & 1.7580054 & 0.9133771 & 0.0000000 &  & \\
4 & 2.3013651 & 2.5166265 & 2.5338848 & 0.0000000 & \\
5 & 0.3835608 & 1.5981702 & 2.0330006 & 1.9178042 & 0.0000000\\
\bottomrule\end{tabular}
\end{minipage}\end{center}

\subsubsection{明氏距离的缺点}
\begin{enumerate}
\item 未能去除量纲的影响。虽然可通过Z变换去除量纲影响，但Z变换相当于对空间按各维度的标准差进行压缩，改变原始变量的空间关系。
\item 未考虑变量间相关性。相关程度高的变量较多，会使计算距离夸大样本间的疏离程度。
\end{enumerate}
\Note{有的资料称“\texttt{dist()}对数据中心化处理”，不确切。\texttt{dist()}直接按输入数据计算距离，既不中心化也不标准化。中心化是各变量减去同一常数，样品间的差值$x_{ik}-x_{jk}$不变，距离也不变；标准化则改变距离，需要事先用\texttt{scale()}完成，如本例。}
标准化后的欧氏距离为
\[
d_{ij}^{*}=\left[\sum_{k=1}^p\frac{(x_{ik}-x_{jk})^2}{s_k^2}\right]^{1/2},
\]
相当于给第$k$个变量加权$1/s_k^2$：方差大的变量被压低，方差小的变量被放大，使每个变量对平方距离的平均贡献相同（对所有样品对平均，每个变量的$(x_{ik}-x_{jk})^2/s_k^2$都等于2）。这种等权是否合理，要看各变量在专业上是否同等重要。
\subsection{马氏距离}
明氏距离特别是欧氏距离的缺点是未考虑指标间相关性。Mahalanobis提出马氏距离。当指标间协方差阵为单位阵时，马氏距离即欧氏距离。
\begin{equation}
d_{ij}^2(M)=(X_i-X_j)\trans\Sigma^{-1}(X_i-X_j).
\end{equation}
马氏距离去除量纲影响的同时，也去除变量间相关性的影响。各坐标差的权重由协方差阵的逆矩阵$\Sigma^{-1}$确定的二次型给出，交叉项也参与加权，不是简单地以某个相关系数的倒数加权。但马氏距离也改变了样本间的空间关系。

要求$\Sigma$正定，$\Sigma^{-1}$和对称平方根$\Sigma^{-1/2}$才存在。
\begin{zm}[马氏距离即白化后的欧氏距离]
令$Z=\Sigma^{-1/2}(X-\mu)$，则$\cov(Z)=\Sigma^{-1/2}\Sigma\Sigma^{-1/2}=I$，各分量方差为1且互不相关。对两个样品，$Z_i-Z_j=\Sigma^{-1/2}(X_i-X_j)$，故
\[
\|Z_i-Z_j\|^2=(X_i-X_j)\trans\Sigma^{-1/2}\Sigma^{-1/2}(X_i-X_j)=(X_i-X_j)\trans\Sigma^{-1}(X_i-X_j)=d_{ij}^2(M).
\]
\end{zm}
实际计算用样本协方差阵$S$代替$\Sigma$。变量完全共线，或样本量不足（$n-1<p$）时，$S$不可逆，马氏距离无法计算；样本量小时即使可逆，$S^{-1}$也很不稳定。本例$n=5$、$p=3$，仅供演示。
\begin{code}
\RComment{\# mahalanobis distances}
n \ROperator{<-} \RFunction{nrow}(da); S \ROperator{<-} \RFunction{cov}(da)
d \ROperator{<-} \RFunction{matrix}(\RKeyword{NA}, n, n)
\RKeyword{for} (i \RKeyword{in} \RNumber{1}\ROperator{:}n) \{
  \RKeyword{for} (j \RKeyword{in} \RNumber{1}\ROperator{:}n) \{
    a \ROperator{<-} \RFunction{as.matrix}(da[i, ] \ROperator{-} da[j, ])
    d[i, j] \ROperator{<-} a \ROperator{\%*\%} \RFunction{solve}(S) \ROperator{\%*\%} \RFunction{t}(a)
  \}
\}
d
\end{code}
\par\addvspace{5pt}\begin{center}\begin{minipage}{\linewidth}\centering
\small
\captionof{table}{平方马氏距离}
\begin{tabular}{lrrrrr}\toprule
 &{\bfseries  1 }&{\bfseries  2 }&{\bfseries  3 }&{\bfseries  4 }&{\bfseries  5}\\\midrule
1 & 0.0000000 & 6.3819087 & 6.0675602 & 7.0421877 & 0.2554311\\
2 & 6.3819087 & 0.0000000 & 7.9860571 & 7.9139379 & 5.7172888\\
3 & 6.0675602 & 7.9860571 & 0.0000000 & 7.8307143 & 6.0601518\\
4 & 7.0421877 & 7.9139379 & 7.8307143 & 0.0000000 & 4.7447623\\
5 & 0.2554311 & 5.7172888 & 6.0601518 & 4.7447623 & 0.0000000\\
\bottomrule\end{tabular}
\end{minipage}\end{center}

\subsection{配合距离}
对于分类变量，尤其是无序分类变量，使用配合距离度量样品间相似度。设两样品都有5个指标值：
\[
S_1=(V,Q,S,T,K),\qquad S_2=(V,M,S,F,K).
\]
二者的配合相似度为
\[
s_{ij}=\frac{m_{ij}}{p}=\frac35=0.6,\qquad\text{配合距离}\ d_{ij}=1-s_{ij}=\frac25=0.4,
\]
其中$m_{ij}$为两样品取值相同的指标数，$p$为指标数；配合距离即取值不同的指标所占比例。
\subsection{二分类变量的样本间距离}
\begin{center}\begin{tabular}{c@{\qquad}cc}\toprule
&{\bfseries Case $j=1$}&{\bfseries Case $j=0$}\\\midrule
Case $i=1$&$(1,1)$：$A$&$(1,0)$：$C$\\
Case $i=0$&$(0,1)$：$B$&$(0,0)$：$D$\\\bottomrule
\end{tabular}\end{center}
简单配合（simple matching）相似度：
\begin{equation}
s_{ij}=\frac1p\sum_{k=1}^p\bigl[L(A)+L(D)\bigr].
\end{equation}
$L(\cdot)$为逻辑变量，取值0或1。

Jaccard配合相似度（R中的\texttt{binary}）：
\begin{equation}
s_{ij}=\frac{\sum L(A)}{\sum L(A)+\sum L(B)+\sum L(C)}
=\frac{\sum L(A)}{p-\sum L(D)}.
\end{equation}
若两样品所有指标都为0（$p=\sum L(D)$），Jaccard系数为$0/0$，没有定义；R的\texttt{dist(..., "binary")}此时把距离记为0（见下表样本9与10）。样本间距离为$d_{ij}=1-s_{ij}$。
\subsubsection{Jaccard距离}
用于不对称的二值变量：1表示具有某种特征（阳性、出现某症状），0表示不具有。两个样品都为0，并不能说明它们相似，例如两人都没有某种罕见症状。Jaccard系数把(0,0)配对排除在外。
\par\addvspace{5pt}\begin{center}\begin{minipage}{\linewidth}\centering
\small
\captionof{table}{十个样本的二值指标}
\begin{tabular}{lrrrrrrrrrr}\toprule
{\bfseries 指标 }&{\bfseries  1 }&{\bfseries  2 }&{\bfseries  3 }&{\bfseries  4 }&{\bfseries  5 }&{\bfseries  6 }&{\bfseries  7 }&{\bfseries  8 }&{\bfseries  9 }&{\bfseries  10}\\\midrule
$x_1$ & 0 & 1 & 0 & 0 & 0 & 1 & 1 & 0 & 0 & 0\\
$x_2$ & 0 & 1 & 0 & 1 & 0 & 0 & 1 & 1 & 0 & 0\\
$x_3$ & 1 & 0 & 1 & 0 & 1 & 1 & 1 & 0 & 0 & 0\\
$x_4$ & 1 & 1 & 1 & 0 & 0 & 0 & 0 & 0 & 0 & 0\\
\bottomrule\end{tabular}
\end{minipage}\end{center}

\begin{code}
db \ROperator{<-} \RFunction{cbind}(x1 \ROperator{=} \RFunction{c}(\RNumber{0,1,0,0,0,1,1,0,0,0}), x2 \ROperator{=} \RFunction{c}(\RNumber{0,1,0,1,0,0,1,1,0,0}),
            x3 \ROperator{=} \RFunction{c}(\RNumber{1,0,1,0,1,1,1,0,0,0}), x4 \ROperator{=} \RFunction{c}(\RNumber{1,1,1,0,0,0,0,0,0,0}))
\RFunction{t}(db)                                \RComment{\# 行为样本，列为指标}
d \ROperator{<-} \RFunction{dist}(db, \RString{"binary"})
\RFunction{round}(d, \RNumber{3})
\end{code}
\par\addvspace{5pt}\begin{center}\begin{minipage}{\linewidth}\centering
\small
\captionof{table}{Jaccard距离}
\begin{tabular}{lrrrrrrrrrr}\toprule
 &{\bfseries  1 }&{\bfseries  2 }&{\bfseries  3 }&{\bfseries  4 }&{\bfseries  5 }&{\bfseries  6 }&{\bfseries  7 }&{\bfseries  8 }&{\bfseries  9 }&{\bfseries  10}\\\midrule
1 & 0.000 &  &  &  &  &  &  &  &  & \\
2 & 0.750 & 0.000 &  &  &  &  &  &  &  & \\
3 & 0.000 & 0.750 & 0.000 &  &  &  &  &  &  & \\
4 & 1.000 & 0.667 & 1.000 & 0.000 &  &  &  &  &  & \\
5 & 0.500 & 1.000 & 0.500 & 1.000 & 0.000 &  &  &  &  & \\
6 & 0.667 & 0.750 & 0.667 & 1.000 & 0.500 & 0.000 &  &  &  & \\
7 & 0.750 & 0.500 & 0.750 & 0.667 & 0.667 & 0.333 & 0.000 &  &  & \\
8 & 1.000 & 0.667 & 1.000 & 0.000 & 1.000 & 1.000 & 0.667 & 0.000 &  & \\
9 & 1.000 & 1.000 & 1.000 & 1.000 & 1.000 & 1.000 & 1.000 & 1.000 & 0.000 & \\
10 & 1.000 & 1.000 & 1.000 & 1.000 & 1.000 & 1.000 & 1.000 & 1.000 & 0.000 & 0.000\\
\bottomrule\end{tabular}
\end{minipage}\end{center}

\subsubsection{简单配合距离}
用于对称的二值变量：0和1两种取值地位相当，如性别、城乡，此时(0,0)与(1,1)同样是相似的证据。选择简单配合还是Jaccard，依据是共同取0能否作为相似的证据，而不是0、1的比例是否均衡。
\begin{code}
n \ROperator{<-} \RFunction{nrow}(db); p \ROperator{<-} \RFunction{ncol}(db)
d \ROperator{<-} \RFunction{matrix}(\RKeyword{NA}, n, n)
\RKeyword{for} (i \RKeyword{in} \RNumber{1}\ROperator{:}n) \{
  \RKeyword{for} (j \RKeyword{in} \RNumber{1}\ROperator{:}n) \{
    L \ROperator{<-} \RFunction{as.numeric}(db[i, ] \ROperator{==} db[j, ])
    d[i, j] \ROperator{<-} \RNumber{1} \ROperator{-} \RFunction{sum}(L) \ROperator{/} p
  \}
\}
d.dist \ROperator{<-} \RFunction{as.dist}(d); d.dist
\end{code}
\par\addvspace{5pt}\begin{center}\begin{minipage}{\linewidth}\centering
\small
\captionof{table}{简单配合距离}
\begin{tabular}{lrrrrrrrrrr}\toprule
 &{\bfseries  1 }&{\bfseries  2 }&{\bfseries  3 }&{\bfseries  4 }&{\bfseries  5 }&{\bfseries  6 }&{\bfseries  7 }&{\bfseries  8 }&{\bfseries  9 }&{\bfseries  10}\\\midrule
1 & 0.00 &  &  &  &  &  &  &  &  & \\
2 & 0.75 & 0.00 &  &  &  &  &  &  &  & \\
3 & 0.00 & 0.75 & 0.00 &  &  &  &  &  &  & \\
4 & 0.75 & 0.50 & 0.75 & 0.00 &  &  &  &  &  & \\
5 & 0.25 & 1.00 & 0.25 & 0.50 & 0.00 &  &  &  &  & \\
6 & 0.50 & 0.75 & 0.50 & 0.75 & 0.25 & 0.00 &  &  &  & \\
7 & 0.75 & 0.50 & 0.75 & 0.50 & 0.50 & 0.25 & 0.00 &  &  & \\
8 & 0.75 & 0.50 & 0.75 & 0.00 & 0.50 & 0.75 & 0.50 & 0.00 &  & \\
9 & 0.50 & 0.75 & 0.50 & 0.25 & 0.25 & 0.50 & 0.75 & 0.25 & 0.00 & \\
10 & 0.50 & 0.75 & 0.50 & 0.25 & 0.25 & 0.50 & 0.75 & 0.25 & 0.00 & 0.00\\
\bottomrule\end{tabular}
\end{minipage}\end{center}

\subsection{无序多分类变量的样本间距离}
简单配合相似度为
\[
s_{ij}=\frac1p\sum_{k=1}^p\kappa(x_{ik},x_{jk}),\qquad
\kappa(x_{ik},x_{jk})=\begin{cases}1,&x_{ik}=x_{jk},\\0,&x_{ik}\ne x_{jk}.\end{cases}
\]
样本间距离$d_{ij}=1-s_{ij}$。
\subsection{有序变量的样本间距离}
先对有序变量作$[0,1]$区间变换。设有序变量$x$取值为$1,2,\ldots,M$，标准化变换为
\[
y=\frac{x-1}{M-1}.
\]
变换后按数值变量计算距离。这种等距映射隐含一个计分约定：相邻等级之间的差异都相等（如“轻—中”与“中—重”的差距相同）。若专业上不能接受，应另行指定各等级的得分。
\section{指标间相似性的度量}
常用夹角余弦、相关系数等。指标$s$与$t$的夹角余弦和Pearson相关系数分别为
\begin{align}
\cos\theta_{st}&=\frac{\sum_{i=1}^n x_{is}x_{it}}{\sqrt{\sum_{i=1}^nx_{is}^2\sum_{i=1}^nx_{it}^2}},\\
r_{st}&=\frac{\sum_{i=1}^n(x_{is}-\bar x_s)(x_{it}-\bar x_t)}
{\sqrt{\sum_{i=1}^n(x_{is}-\bar x_s)^2\sum_{i=1}^n(x_{it}-\bar x_t)^2}}.
\end{align}
\Note{夹角余弦使用原始指标，Pearson相关使用中心化后的指标。思考二者的关系。}
\begin{zm}[中心化向量的夹角余弦等于Pearson相关系数]
把$x_{is}$换成$x_{is}-\bar x_s$、$x_{it}$换成$x_{it}-\bar x_t$代入夹角余弦公式，所得正是$r_{st}$。
\end{zm}
原始指标都取正值且均数远大于变异时，夹角余弦主要反映均数，几乎都接近1（如下例三个指标两两在0.987以上），不能反映变量间的共变；中心化后的相关系数才能。

变量间距离：$c_{st}=r$，$d_{st}=\sqrt{1-c_{st}^2}$。三种常用的相关型距离对负相关的处理不同：
\begin{center}\small\begin{tabular}{lccc}\toprule
{\bfseries 距离}&$r=0.9$&$r=0$&$r=-0.9$\\\midrule
$\sqrt{1-r^2}$&0.436&1&0.436\\
$1-|r|$&0.1&1&0.1\\
$1-r$&0.1&1&1.9\\\bottomrule
\end{tabular}\end{center}
前两种把强负相关与强正相关同样看作“相近”，适合以信息冗余为目的的变量聚类，例如“往返跑时间”越短越好，与其他体能指标负相关，但反映的是同一种能力。$1-r$则把强负相关看作最远，适合方向本身有意义的情形。下例用$\sqrt{1-r^2}$；服装标准例用$1-|r|$（其相关系数全为正，与$1-r$结果相同）；体质数据练习用$1-r$：往返跑、反应时等计时指标与纵跳、肺活量、俯卧撑负相关（如往返跑与纵跳$r=-0.36$），用$1-r$时距离都大于1，会被分开。若关心的是信息冗余，应改用$1-|r|$。
\begin{code}
a \ROperator{<-} \RFunction{as.matrix}(da); p \ROperator{<-} \RFunction{ncol}(a); s \ROperator{<-} \RFunction{matrix}(\RKeyword{NA}, p, p)
\RKeyword{for} (i \RKeyword{in} \RNumber{1}\ROperator{:}p) \{
  \RKeyword{for} (j \RKeyword{in} \RNumber{1}\ROperator{:}p) \{
    s[i, j] \ROperator{<-} a[, i] \ROperator{\%*\%} a[, j] \ROperator{/}
      \RFunction{sqrt}((a[, i] \ROperator{\%*\%} a[, i]) \ROperator{*} (a[, j] \ROperator{\%*\%} a[, j]))
  \}
\}
\RFunction{colnames}(s) \ROperator{<-} \RFunction{colnames}(da); \RFunction{rownames}(s) \ROperator{<-} \RFunction{colnames}(da)
s
d \ROperator{<-} \RFunction{sqrt}(\RNumber{1} \ROperator{-} s\ROperator{\textasciicircum{}}\RNumber{2})
d.dist \ROperator{<-} \RFunction{as.dist}(d); d.dist
\RComment{\# 变量的距离，Pearson r}
d \ROperator{<-} \RFunction{sqrt}(\RNumber{1} \ROperator{-} \RFunction{cor}(da)\ROperator{\textasciicircum{}}\RNumber{2})
d.dist \ROperator{<-} \RFunction{as.dist}(d); d.dist
\end{code}
\par\addvspace{5pt}\begin{center}\begin{minipage}{\linewidth}\centering
\small
\captionof{table}{夹角余弦相似系数}
\begin{tabular}{lrrr}\toprule
 &{\bfseries  tz }&{\bfseries  xw }&{\bfseries  wfp}\\\midrule
tz & 1.0000000 & 0.9963195 & 0.9869652\\
xw & 0.9963195 & 1.0000000 & 0.9918570\\
wfp & 0.9869652 & 0.9918570 & 1.0000000\\
\bottomrule\end{tabular}
\end{minipage}\end{center}
\par\addvspace{5pt}\begin{center}\begin{minipage}{\linewidth}\centering
\small
\captionof{table}{夹角余弦对应的距离}
\begin{tabular}{lrrr}\toprule
 &{\bfseries  tz }&{\bfseries  xw }&{\bfseries  wfp}\\\midrule
tz & 0.00000000 &  & \\
xw & 0.08571709 & 0.00000000 & \\
wfp & 0.16093352 & 0.12735628 & 0.00000000\\
\bottomrule\end{tabular}
\end{minipage}\end{center}
\par\addvspace{5pt}\begin{center}\begin{minipage}{\linewidth}\centering
\small
\captionof{table}{Pearson相关对应的距离}
\begin{tabular}{lrrr}\toprule
 &{\bfseries  tz }&{\bfseries  xw }&{\bfseries  wfp}\\\midrule
tz & 0.0000000 &  & \\
xw & 0.4799079 & 0.0000000 & \\
wfp & 0.9932792 & 0.9981048 & 0.0000000\\
\bottomrule\end{tabular}
\end{minipage}\end{center}

对等级资料可用Spearman相关系数，对分类资料可用列联系数
\[
C=\sqrt{\frac{\chi^2}{\chi^2+n}}
\]
表示相似程度。
\section{层次聚类法}
层次聚类即系统聚类，是最常用的样品或变量聚类方法。
\begin{enumerate}
\item 构造$n$个类，每类仅包含一个样品或指标。
\item 计算$n$个类两两间的距离，构成距离矩阵。
\item 合并距离最近的两类为一新类，此时剩$n-1$类。
\item 计算新类与当前各类的距离，重复上一步，直至只剩一类。
\end{enumerate}
\subsection{类间距离的计算方法}
每类含一个样品时，类间距离即样品间距离。否则采用最短距离法、最长距离法、类平均法、重心法、Ward法等。
\begin{figure}[!htbp]\centering
\begin{minipage}{.48\textwidth}\centering \begin{tikzpicture}[x=.75cm,y=.75cm,every node/.style={font=\small}]\draw[Muted,line width=.6pt] (0,0) circle (1);\draw[Muted,line width=.6pt] (4.6,0) circle (1);\fill[Brick] (-0.2,0.5) circle (2.3pt) node[above right] {$S_1$};\fill[Brick] (0.5,-0.05) circle (2.3pt) node[above right] {$S_2$};\fill[Brick] (-0.5,-0.6) circle (2.3pt) node[above right] {$S_3$};\fill[Forest] (4.25,0.5) circle (2.3pt) node[above right] {$S_4$};\fill[Forest] (5.05,-0.2) circle (2.3pt) node[above right] {$S_5$};\node at (0,-1.35) {$G_1$};\node at (4.6,-1.35) {$G_2$};\draw[Plum,line width=1.1pt] (.5,-.05)--(4.25,.5);\end{tikzpicture}\\[2pt]最短距离\end{minipage}\hfill\begin{minipage}{.48\textwidth}\centering \begin{tikzpicture}[x=.75cm,y=.75cm,every node/.style={font=\small}]\draw[Muted,line width=.6pt] (0,0) circle (1);\draw[Muted,line width=.6pt] (4.6,0) circle (1);\fill[Brick] (-0.2,0.5) circle (2.3pt) node[above right] {$S_1$};\fill[Brick] (0.5,-0.05) circle (2.3pt) node[above right] {$S_2$};\fill[Brick] (-0.5,-0.6) circle (2.3pt) node[above right] {$S_3$};\fill[Forest] (4.25,0.5) circle (2.3pt) node[above right] {$S_4$};\fill[Forest] (5.05,-0.2) circle (2.3pt) node[above right] {$S_5$};\node at (0,-1.35) {$G_1$};\node at (4.6,-1.35) {$G_2$};\draw[Plum,line width=1.1pt] (-.5,-.6)--(5.05,-.2);\end{tikzpicture}\\[2pt]最长距离\end{minipage}\\[6mm]\begin{minipage}{.48\textwidth}\centering \begin{tikzpicture}[x=.75cm,y=.75cm,every node/.style={font=\small}]\draw[Muted,line width=.6pt] (0,0) circle (1);\draw[Muted,line width=.6pt] (4.6,0) circle (1);\fill[Brick] (-0.2,0.5) circle (2.3pt) node[above right] {$S_1$};\fill[Brick] (0.5,-0.05) circle (2.3pt) node[above right] {$S_2$};\fill[Brick] (-0.5,-0.6) circle (2.3pt) node[above right] {$S_3$};\fill[Forest] (4.25,0.5) circle (2.3pt) node[above right] {$S_4$};\fill[Forest] (5.05,-0.2) circle (2.3pt) node[above right] {$S_5$};\node at (0,-1.35) {$G_1$};\node at (4.6,-1.35) {$G_2$};\fill[Brick] (-.0667,-.05) circle (2.3pt) node[below left] {$C_1$};\fill[Forest] (4.65,.15) circle (2.3pt) node[above] {$C_2$};\draw[Plum,line width=1.1pt] (-.0667,-.05)--(4.65,.15);\draw[Muted,densely dotted,line width=.55pt] (-0.2,0.5)--(-.0667,-.05);\draw[Muted,densely dotted,line width=.55pt] (0.5,-0.05)--(-.0667,-.05);\draw[Muted,densely dotted,line width=.55pt] (-0.5,-0.6)--(-.0667,-.05);\draw[Muted,densely dotted,line width=.55pt] (4.25,0.5)--(4.65,.15);\draw[Muted,densely dotted,line width=.55pt] (5.05,-0.2)--(4.65,.15);\end{tikzpicture}\\[2pt]重心距离\end{minipage}\hfill\begin{minipage}{.48\textwidth}\centering \begin{tikzpicture}[x=.75cm,y=.75cm,every node/.style={font=\small}]\draw[Muted,line width=.6pt] (0,0) circle (1);\draw[Muted,line width=.6pt] (4.6,0) circle (1);\fill[Brick] (-0.2,0.5) circle (2.3pt) node[above right] {$S_1$};\fill[Brick] (0.5,-0.05) circle (2.3pt) node[above right] {$S_2$};\fill[Brick] (-0.5,-0.6) circle (2.3pt) node[above right] {$S_3$};\fill[Forest] (4.25,0.5) circle (2.3pt) node[above right] {$S_4$};\fill[Forest] (5.05,-0.2) circle (2.3pt) node[above right] {$S_5$};\node at (0,-1.35) {$G_1$};\node at (4.6,-1.35) {$G_2$};\draw[Muted,densely dotted,line width=.55pt] (-0.2,0.5)--(4.25,0.5);\draw[Muted,densely dotted,line width=.55pt] (-0.2,0.5)--(5.05,-0.2);\draw[Muted,densely dotted,line width=.55pt] (0.5,-0.05)--(4.25,0.5);\draw[Muted,densely dotted,line width=.55pt] (0.5,-0.05)--(5.05,-0.2);\draw[Muted,densely dotted,line width=.55pt] (-0.5,-0.6)--(4.25,0.5);\draw[Muted,densely dotted,line width=.55pt] (-0.5,-0.6)--(5.05,-0.2);\end{tikzpicture}\\[2pt]类平均距离\end{minipage}
\caption{类间距离的四种计算方法}
\end{figure}\FloatBarrier
类平均法对输入的相异度取所有跨类样品对的平均：
\[
D_{12}=\frac1{n_1n_2}\sum_{i\in G_1}\sum_{j\in G_2}d_{ij}.
\]
若输入的是平方欧氏距离，则得到常见的平方距离平均式（本图$n_1=3$、$n_2=2$）：
\[D_{12}^2=\frac16(d_{14}^2+d_{15}^2+d_{24}^2+d_{25}^2+d_{34}^2+d_{35}^2).\]
两种输入的合并顺序一般不同，应在报告中写明。

\begin{description}[style=nextline,leftmargin=0pt,labelwidth=0pt,labelsep=0pt]
\item[类平均法（average）]平均距离最小的类被合并。利用所有跨类样品对的信息，性质介于最短与最长距离法之间。
\item[重心法（centroid）]两类间距离为两类均值向量（重心）之间的距离。新类重心按样品数加权，大类与小类合并后重心几乎不动；合并高度可能不单调，聚类树出现“倒置”。用递推公式计算时，输入应为平方欧氏距离。
\item[最短距离法（single）]两类间距离取跨类样品对距离的最小值。只要有一对样品相近，两类就合并，容易产生链式效应，把细长或首尾相连的点串成一类；能识别非椭圆形的类，但对两类之间的“桥接”点敏感。
\item[最长距离法（complete）]两类间距离取跨类样品对距离的最大值。合并高度就是新类的类内最大差异（直径），倾向于得到紧凑、直径相近的类，对离群值敏感。（在Lance--Williams的术语中，最短距离法是“空间压缩”的，最长距离法是“空间扩张”的。）
\item[Ward法]每一步合并使类内离差平方和$W$增加最少的两类。合并$G_a$、$G_b$的增量为
\[
\Delta W=\frac{n_an_b}{n_a+n_b}\|\bar X_a-\bar X_b\|^2.
\]
系数$n_an_b/(n_a+n_b)$在一类样品很少时较小，因此小类倾向于先被合并，各类样品数较均衡，形状接近球形。分类效率高，对极端值敏感。R中\texttt{hclust(dist(x), method = \string"ward.D2\string")}以未平方的欧氏距离为输入实现该准则；\texttt{\string"ward.D\string"}则须输入平方欧氏距离\texttt{dist(x)\string^2}才与之对应。
\end{description}
\begin{zm}[Ward合并增量]
记$n=n_a+n_b$，合并后重心$\bar X=(n_a\bar X_a+n_b\bar X_b)/n$。由平方和分解，
\[
W_{a\cup b}=\sum_{i\in G_a\cup G_b}\|X_i-\bar X\|^2=W_a+W_b+n_a\|\bar X_a-\bar X\|^2+n_b\|\bar X_b-\bar X\|^2,
\]
其中交叉项$\sum_{i\in G_a}(X_i-\bar X_a)=0$而消失。又$\bar X_a-\bar X=\frac{n_b}{n}(\bar X_a-\bar X_b)$、$\bar X_b-\bar X=\frac{n_a}{n}(\bar X_b-\bar X_a)$，故
\[
\Delta W=W_{a\cup b}-W_a-W_b=\frac{n_an_b^2+n_bn_a^2}{n^2}\|\bar X_a-\bar X_b\|^2=\frac{n_an_b}{n_a+n_b}\|\bar X_a-\bar X_b\|^2.
\]
\end{zm}
\begin{sikao}[类间距离定义对聚类结果的影响]
对前文体重、胸围、往返跑的5人资料（标准化后的欧氏距离），分别用四种方法作系统聚类：
\par\noindent\begin{rcode}
x <- scale(da); d <- dist(x)
for(m in c("single","complete","average","ward.D2")){
  h <- hclust(d,method=m)
  cat(m,": 合并高度",round(h$height,3)," | 分2类:",cutree(h,2),"\n")
}
\end{rcode}
\begin{routput}
single : 合并高度 0.583 1.081 2.265 2.409  | 分2类: 1 2 2 1 1 
complete : 合并高度 0.583 1.081 2.576 3.487  | 分2类: 1 1 1 2 1 
average : 合并高度 0.583 1.081 2.457 2.896  | 分2类: 1 1 1 2 1 
ward.D2 : 合并高度 0.583 1.081 2.936 3.725  | 分2类: 1 2 2 1 1 
\end{routput}
四种方法的前两步完全相同：先合并距离最近的1号与5号（0.583），再合并2号与3号（1.081），因为每类仅含一个样品时，类间距离即样品间距离。自第三步起结果出现分歧：最短距离法与Ward法分为\{1,4,5\}与\{2,3\}两类，最长距离法与类平均法分为\{1,2,3,5\}与\{4\}两类。

仅5人、3个变量的资料，分类结果即随方法而变，样本较大时更是如此。实际分析中可比较多种方法的结果：主要类别在不同方法下稳定出现时，结果较为可信；结果随方法变化较大时，表明数据中的类别结构不清晰。
\end{sikao}
\subsection{类的个数}
给定阈值$T$，在聚类树上“切一刀”，较主观。也可观察样品散点图，但仅适合2至3个变量的情况。更客观的辅助指标（轮廓系数、重抽样稳定性）见\ref{sec:k-number}节。
\subsection{系统聚类实例}
\Example{例1\quad Q型系统聚类：细菌分类}
用气相色谱法分析细菌全细胞脂肪酸含量，研究细菌分类与鉴定。24株菌株包括空肠弯曲菌8株（CJ1至CJ8）、结肠弯曲菌3株（CC1至CC3）、幽门螺旋菌9株（HP1至HP9）及其他肠道杆菌4株（XX1至XX4）。测得12种脂肪酸百分含量（$X_1$至$X_{12}$），要求以此对24株菌株进行聚类分析。数据：\texttt{bacteria.sav}。
\clearpage\begin{landscape}\thispagestyle{plain}
\Example{例9.2\quad R型系统聚类：服装标准}
制定我国服装标准时，测量3454名成年女子的14个部位。根据相关系数矩阵，对变量进行聚类分析。采用距离
\[
d_{ij}=1-|r_{ij}|.
\]

\par\addvspace{5pt}\begin{center}\begin{minipage}{\linewidth}\centering
\small
\captionof{table}{3454名成年女子14项身体测量指标的相关矩阵}
\begin{tabular}{lrrrrrrrrrrrrrr}\toprule
 &{\bfseries  上体长 }&{\bfseries  手臂长 }&{\bfseries  胸围 }&{\bfseries  颈围 }&{\bfseries  总肩宽 }&{\bfseries  前胸宽 }&{\bfseries  后背宽 }&{\bfseries  前腰节高 }&{\bfseries  后腰节高 }&{\bfseries  总体高 }&{\bfseries  身高 }&{\bfseries  下体长 }&{\bfseries  腰围 }&{\bfseries  臀围}\\\midrule
上体长 & 1.000 &  &  &  &  &  &  &  &  &  &  &  &  & \\
手臂长 & 0.366 & 1.000 &  &  &  &  &  &  &  &  &  &  &  & \\
胸围 & 0.242 & 0.233 & 1.000 &  &  &  &  &  &  &  &  &  &  & \\
颈围 & 0.280 & 0.194 & 0.590 & 1.000 &  &  &  &  &  &  &  &  &  & \\
总肩宽 & 0.360 & 0.324 & 0.476 & 0.435 & 1.000 &  &  &  &  &  &  &  &  & \\
前胸宽 & 0.282 & 0.263 & 0.483 & 0.470 & 0.452 & 1.000 &  &  &  &  &  &  &  & \\
后背宽 & 0.245 & 0.265 & 0.540 & 0.478 & 0.535 & 0.663 & 1.000 &  &  &  &  &  &  & \\
前腰节高 & 0.448 & 0.345 & 0.452 & 0.404 & 0.431 & 0.322 & 0.266 & 1.000 &  &  &  &  &  & \\
后腰节高 & 0.486 & 0.367 & 0.365 & 0.357 & 0.429 & 0.283 & 0.287 & 0.820 & 1.000 &  &  &  &  & \\
总体高 & 0.648 & 0.662 & 0.216 & 0.316 & 0.429 & 0.283 & 0.263 & 0.527 & 0.547 & 1.000 &  &  &  & \\
身高 & 0.679 & 0.681 & 0.243 & 0.313 & 0.430 & 0.302 & 0.294 & 0.520 & 0.558 & 0.957 & 1.000 &  &  & \\
下体长 & 0.486 & 0.636 & 0.174 & 0.243 & 0.375 & 0.290 & 0.255 & 0.403 & 0.417 & 0.857 & 0.852 & 1.000 &  & \\
腰围 & 0.133 & 0.153 & 0.732 & 0.477 & 0.339 & 0.392 & 0.446 & 0.266 & 0.241 & 0.054 & 0.099 & 0.055 & 1.000 & \\
臀围 & 0.376 & 0.252 & 0.676 & 0.581 & 0.441 & 0.447 & 0.440 & 0.424 & 0.372 & 0.363 & 0.376 & 0.321 & 0.627 & 1.000\\
\bottomrule\end{tabular}
\end{minipage}\end{center}
\Source{资料来源：张尧庭、方开泰（1982），《多元统计分析引论》，北京：科学出版社，116页。}
\end{landscape}
\Example{裁判评分的相似度}
SPSS自带数据集\texttt{judge.sav}记录中、美、法等七国裁判及未经严格训练的体育爱好者对选手的评分。根据评分差异，将其分为适当的若干类。
\Example{练习\quad 体质数据}
使用\texttt{体质数据.sav}，对\texttt{x1}至\texttt{x12}作R型聚类。
\begin{code}
\RComment{\# 数据：体质数据.sav。课程版本第1列即脉搏mb，da[, -1]会把它去掉，应按所用版本核对列名}
rm \ROperator{<-} \RFunction{cor}(da[, \ROperator{-}\RNumber{1}]); rm
d \ROperator{<-} \RFunction{as.dist}(\RNumber{1} \ROperator{-} rm); d        \RComment{\# 若关心信息冗余，改用 1 - abs(rm)}
cv \ROperator{<-} \RFunction{hclust}(d); cv
\RFunction{plot}(cv, hang \ROperator{=} \ROperator{-}\RNumber{1})
\end{code}
\par\addvspace{5pt}\begin{center}\begin{minipage}{\linewidth}\centering
\small
\captionof{table}{体质数据的变量信息}
\begin{tabular}{llll}\toprule
{\bfseries 变量 }&{\bfseries  指标 }&{\bfseries  变量 }&{\bfseries  指标}\\\midrule
\texttt{mb} & 脉搏 & \texttt{tjzs} & 台阶指数\\
\texttt{ssy} & 收缩压 & \texttt{wl} & 握力\\
\texttt{szy} & 舒张压 & \texttt{bl} & 背力\\
\texttt{fhl} & 肺活量 & \texttt{tqq} & 体前屈\\
\texttt{sg} & 身高 & \texttt{zt} & 纵跳\\
\texttt{tz} & 体重 & \texttt{djzl} & 单脚站立\\
\texttt{xw} & 胸围 & \texttt{wfp} & 往返跑\\
\texttt{yw} & 腰围 & \texttt{fys} & 反应时1\\
\texttt{tw} & 臀围 & \texttt{fwc} & 俯卧撑\\
\bottomrule\end{tabular}
\end{minipage}\end{center}

\subsection{变量聚类用于变量优选}
变量聚类是数据降维的简便方法，直接去掉某些变量。代表性指标的选择考虑：专业意义、易测性与经济性、统计代表性。
\begin{equation}
R_i^2(J)=\frac{\displaystyle\sum_{\substack{k\in J\\k\ne i}}r_{ik}^2}{M_J-1}.
\end{equation}
第$J$类有$M_J$个变量（$M_J>1$），$R_i^2(J)$是第$i$个变量与同类其余$M_J-1$个变量相关系数平方的平均值。它越大，该变量越能代表本类，可在专业意义、易测性相近时选$R_i^2(J)$最大者作为代表性指标。
\subsection{计算问题}
聚类分析的计算量与样本量的平方或立方成正比，系统聚类不太适合大样本。快速聚类法需事先指定类的个数。
\section{大样本快速聚类}
K均值聚类（K-means clustering）的分类思想：
\begin{enumerate}
\item 事先指定分类数$k$与初始点。
\item 样本依距离（欧氏距离平方）关系归入$k$类。
\item 以类均数向量为中心重新计算。
\item 重复上一步迭代，直到分类不变。
\end{enumerate}
分类数$k$按专业需要或数据探索确定。
\subsection{目标函数与迭代}
记分组为$C_1,\ldots,C_k$，第$l$类样品数为$n_l$、中心为$c_l$。K均值要最小化类内平方和
\[
W=\sum_{l=1}^k\sum_{i\in C_l}\|x_i-c_l\|^2.
\]
\begin{zm}[固定分组时，均值使$W$最小]
对任一类$C_l$和任一向量$c$，
\[
\sum_{i\in C_l}\|x_i-c\|^2=\sum_{i\in C_l}\|x_i-\bar x_l\|^2+n_l\|\bar x_l-c\|^2,
\]
交叉项因$\sum_{i\in C_l}(x_i-\bar x_l)=0$而消失。右端在$c=\bar x_l$时最小。
\end{zm}
固定各类中心时，把每个样品分到最近（欧氏距离平方最小）的中心，使$W$的每一项都最小。K均值交替进行“分配”和“更新中心”两步，每一步都不增大$W$。不同的分组只有有限多种，算法在有限步内停止，但停在的是局部最优，不一定是全局最优，结果依赖初始点。

停止条件与优化目标是两回事。常用的停止条件是：分组不再改变（或中心的变化小于容许误差），或达到最大迭代次数（R中由\texttt{iter.max}控制，默认10）。R的\texttt{kmeans()}默认用Hartigan--Wong算法，按样品逐个调整，同样使$W$单调不增。
\begin{zm}[$T=W+B$]
记总均值为$\bar x$，$T=\sum_i\|x_i-\bar x\|^2$，$B=\sum_ln_l\|\bar x_l-\bar x\|^2$。则
\[
T=\sum_l\sum_{i\in C_l}\|(x_i-\bar x_l)+(\bar x_l-\bar x)\|^2=W+B+2\sum_l(\bar x_l-\bar x)\trans\sum_{i\in C_l}(x_i-\bar x_l)=W+B.
\]
\end{zm}
\subsection{初始点的指定}
初始点（initial cluster centres）是迭代开始时各类别的重心。不同的初始点可能导致不同结果。可以由程序随机指定，也可人工选择$k$个典型个案、按频数分布划分区间、采用核密度法等。
\Example{耳朵测量数据}
某医院康复专科门诊为修复耳缺损，测量300例患者正常侧耳朵的五项指标：耳长（EC）、耳宽（EK）、耳外展距（EZ）、耳型（EX）及耳垂型（ECX）。另计算两个指数：
\[
\mathrm{EI}=\frac{\text{耳宽}}{\text{耳长}}\times100\%,\qquad
\mathrm{AI}=\frac{\text{耳外展距}}{\text{耳宽}}\times100\%.
\]
要求以EC、EK、EZ、EI、AI五个变量为依据进行聚类，找出各类具有代表性的耳朵，供临床修复耳缺损参考。数据：\texttt{ear.sav}（300例）。
\Note{原例正文写作“EC、EX、EZ、EI、AI”，但结果表各列为EC、EK、EZ、EI、AI。按数据文件的变量标签（EC耳长、EK耳宽、EZ耳外展距、EI耳指数、AI外展指数、EX耳型、ECX耳垂型），并用课程数据复算，下列全部结果都由EC、EK、EZ、EI、AI得到，正文的EX应为EK。EX、ECX是分类变量，不宜直接放入K均值。}
EI、AI由EC、EK、EZ算出，并非独立的测量。标准化后每个变量在欧氏距离中权重相同，耳宽EK同时出现在EK、EI、AI三个变量里，“宽窄”和“外展”方面的形状差异实际上被加了权。若希望长、宽、外展距等权，可只用三个原始测量，或只用EC与两个指数；结果应结合聚类的目的来选择。
\subsection{数据探索}
\begin{code}
da \ROperator{<-} foreign\ROperator{::}\RFunction{read.spss}(\RString{"ear.sav"}, to.data.frame \ROperator{=} \RKeyword{TRUE})
v \ROperator{<-} \RFunction{c}(\RString{"EC"}, \RString{"EK"}, \RString{"EZ"}, \RString{"EI"}, \RString{"AI"})   \RComment{\# 按变量名选取，不同版本文件列序不同}
\RComment{\# 检查方差，看是否需要标化}
\RFunction{cov}(da[, v])
db \ROperator{<-} \RFunction{scale}(da[, v])
\end{code}
\par\addvspace{5pt}\begin{center}\begin{minipage}{\linewidth}\centering
\small
\captionof{table}{耳测量数据的协方差阵}
\begin{tabular}{lrrrrr}\toprule
 &{\bfseries  EC }&{\bfseries  EK }&{\bfseries  EZ }&{\bfseries  EI }&{\bfseries  AI}\\\midrule
EC & 0.28571639 & 0.07537258 & 0.05918261 & -1.10094177 & 0.42935444\\
EK & 0.07537258 & 0.07985106 & 0.01307202 & 0.64979547 & -1.14300482\\
EZ & 0.05918261 & 0.01307202 & 0.08433935 & -0.26581275 & 2.37003609\\
EI & -1.10094177 & 0.64979547 & -0.26581275 & 19.15508520 & -21.41800780\\
AI & 0.42935444 & -1.14300482 & 2.37003609 & -21.41800780 & 96.96404680\\
\bottomrule\end{tabular}
\end{minipage}\end{center}

\Note{方差差异巨大（EI、AI以\%计，方差为19和97；EC、EK、EZ以cm计，方差不足0.3），需要标化，否则距离几乎只由AI决定。}
\subsection{按百分位数指定初始点}
\begin{code}
c1 \ROperator{<-} \RFunction{round}(\RFunction{quantile}(db[, \RNumber{1}], \RFunction{c}(\RNumber{0.1}, \RNumber{0.25}, \RNumber{0.5}, \RNumber{0.75}, \RNumber{0.9})), \RNumber{1})
c2 \ROperator{<-} \RFunction{round}(\RFunction{quantile}(db[, \RNumber{2}], \RFunction{c}(\RNumber{0.1}, \RNumber{0.25}, \RNumber{0.5}, \RNumber{0.75}, \RNumber{0.9})), \RNumber{1})
c3 \ROperator{<-} \RFunction{round}(\RFunction{quantile}(db[, \RNumber{3}], \RFunction{c}(\RNumber{0.1}, \RNumber{0.25}, \RNumber{0.5}, \RNumber{0.75}, \RNumber{0.9})), \RNumber{1})
c4 \ROperator{<-} \RFunction{round}(\RFunction{quantile}(db[, \RNumber{4}], \RFunction{c}(\RNumber{0.1}, \RNumber{0.25}, \RNumber{0.5}, \RNumber{0.75}, \RNumber{0.9})), \RNumber{1})
c5 \ROperator{<-} \RFunction{round}(\RFunction{quantile}(db[, \RNumber{5}], \RFunction{c}(\RNumber{0.1}, \RNumber{0.25}, \RNumber{0.5}, \RNumber{0.75}, \RNumber{0.9})), \RNumber{1})
ca \ROperator{<-} \RFunction{cbind}(c1, c2, c3, c4, c5); ca
\RFunction{kmeans}(db, ca)
\end{code}
\par\addvspace{5pt}\begin{center}\begin{minipage}{\linewidth}\centering
\small
\captionof{table}{按百分位数指定的初始中心}
\begin{tabular}{lrrrrr}\toprule
{\bfseries 百分位数 }&{\bfseries  c1 }&{\bfseries  c2 }&{\bfseries  c3 }&{\bfseries  c4 }&{\bfseries  c5}\\\midrule
10\% & -1.1 & -1.1 & -1.3 & -1.3 & -1.2\\
25\% & -0.6 & -0.4 & -0.7 & -0.6 & -0.8\\
50\% & 0.0 & -0.1 & 0.0 & 0.1 & 0.0\\
75\% & 0.7 & 0.6 & 0.7 & 0.7 & 0.7\\
90\% & 1.1 & 1.0 & 1.1 & 1.1 & 1.2\\
\bottomrule\end{tabular}
\end{minipage}\end{center}

\subsection{分类结果}
\par\addvspace{5pt}\begin{center}\begin{minipage}{\linewidth}\centering
\small
\captionof{table}{五类的样本数与标化中心}
\begin{tabular}{lrrrrrr}\toprule
{\bfseries 类别 }&{\bfseries  例数 }&{\bfseries  EC }&{\bfseries  EK }&{\bfseries  EZ }&{\bfseries  EI }&{\bfseries  AI}\\\midrule
1 & 24 & -0.68160270 & -1.89445350 & -0.02754703 & -1.34450640 & 1.16177120\\
2 & 87 & -0.37033670 & 0.08444380 & -1.13378211 & 0.45323760 & -1.10602880\\
3 & 73 & -0.74930260 & -0.31329150 & 0.26097356 & 0.42474140 & 0.38155090\\
4 & 50 & 0.76329530 & 1.37778440 & 0.48436869 & 0.60310020 & -0.27338280\\
5 & 66 & 0.98654670 & -0.11967680 & 0.84894952 & -1.03522200 & 0.82057460\\
\bottomrule\end{tabular}
\end{minipage}\end{center}

\par\addvspace{5pt}\begin{center}\begin{minipage}{\linewidth}\centering
\small
\captionof{table}{300例样本的类别（按观测顺序）}
\begin{tabular}{lrrrrrrrrrrrrrrr}\toprule
{\bfseries 编号 }&{\bfseries  1 }&{\bfseries  2 }&{\bfseries  3 }&{\bfseries  4 }&{\bfseries  5 }&{\bfseries  6 }&{\bfseries  7 }&{\bfseries  8 }&{\bfseries  9 }&{\bfseries  10 }&{\bfseries  11 }&{\bfseries  12 }&{\bfseries  13 }&{\bfseries  14 }&{\bfseries  15}\\\midrule
1--15 & 4 & 3 & 3 & 5 & 5 & 3 & 3 & 2 & 3 & 4 & 2 & 1 & 3 & 5 & 5\\
16--30 & 3 & 3 & 5 & 2 & 5 & 3 & 5 & 4 & 2 & 3 & 3 & 1 & 5 & 1 & 3\\
31--45 & 4 & 1 & 3 & 5 & 1 & 5 & 3 & 3 & 3 & 4 & 2 & 5 & 1 & 3 & 5\\
46--60 & 5 & 1 & 1 & 4 & 1 & 5 & 5 & 5 & 1 & 5 & 3 & 3 & 5 & 2 & 3\\
61--75 & 4 & 2 & 5 & 2 & 4 & 5 & 2 & 1 & 1 & 2 & 2 & 3 & 3 & 5 & 4\\
76--90 & 2 & 2 & 3 & 3 & 5 & 1 & 2 & 3 & 3 & 3 & 5 & 2 & 2 & 3 & 2\\
91--105 & 3 & 2 & 4 & 3 & 2 & 5 & 5 & 4 & 1 & 3 & 3 & 4 & 5 & 4 & 4\\
106--120 & 2 & 2 & 4 & 2 & 5 & 3 & 5 & 3 & 4 & 2 & 4 & 3 & 3 & 5 & 2\\
121--135 & 3 & 4 & 2 & 4 & 5 & 3 & 5 & 2 & 2 & 4 & 5 & 4 & 2 & 2 & 2\\
136--150 & 2 & 3 & 2 & 3 & 4 & 4 & 3 & 4 & 3 & 5 & 4 & 3 & 3 & 3 & 4\\
151--165 & 2 & 4 & 5 & 2 & 3 & 5 & 3 & 3 & 2 & 5 & 5 & 2 & 2 & 3 & 3\\
166--180 & 3 & 4 & 2 & 2 & 2 & 2 & 3 & 2 & 3 & 2 & 2 & 5 & 1 & 2 & 3\\
181--195 & 2 & 4 & 2 & 5 & 1 & 3 & 2 & 5 & 5 & 4 & 4 & 2 & 2 & 2 & 4\\
196--210 & 4 & 4 & 3 & 5 & 3 & 5 & 5 & 5 & 2 & 2 & 4 & 4 & 3 & 5 & 2\\
211--225 & 2 & 2 & 2 & 3 & 4 & 2 & 5 & 2 & 3 & 2 & 2 & 5 & 2 & 4 & 4\\
226--240 & 2 & 4 & 2 & 2 & 2 & 3 & 3 & 3 & 5 & 4 & 3 & 5 & 2 & 1 & 5\\
241--255 & 4 & 2 & 5 & 2 & 5 & 3 & 5 & 2 & 5 & 5 & 2 & 2 & 1 & 5 & 5\\
256--270 & 2 & 5 & 2 & 3 & 2 & 5 & 2 & 4 & 5 & 1 & 2 & 1 & 3 & 4 & 2\\
271--285 & 2 & 2 & 1 & 5 & 2 & 1 & 5 & 3 & 2 & 5 & 2 & 4 & 4 & 3 & 4\\
286--300 & 4 & 4 & 3 & 3 & 5 & 2 & 3 & 1 & 3 & 5 & 1 & 4 & 2 & 2 & 2\\
\bottomrule\end{tabular}
\end{minipage}\end{center}

\subsection{优化目标与聚类效果评估}
\begin{description}[style=nextline,leftmargin=0pt,labelwidth=0pt,labelsep=0pt]
\item[Within cluster sum of squares]类内离均差平方和，等同于类内各样本点到类重心（类别中心）的欧氏距离平方和。
\item[Total SS]总平方和，同方差分析中的定义。
\item[Between SS]类间离差平方和。
\end{description}
优化目标（不是收敛准则）：K均值最小化$W$。总平方和$T$与分组无关，由$T=W+B$，等价于
\[
\min\left(\frac{\text{within\_SS}}{\text{total\_SS}}\right)
\quad\text{或}\quad
\max\left(\frac{\text{between\_SS}}{\text{total\_SS}}\right).
\]
何时停止迭代由前述停止条件决定。
\begin{center}\begin{tabular}{cr}\toprule
{\bfseries 类别}&{\bfseries 类内平方和}\\\midrule
1&87.57103\\2&171.28537\\3&110.39236\\4&113.89450\\5&127.54928\\\bottomrule
\end{tabular}\end{center}
\[
\frac{\text{between\_SS}}{\text{total\_SS}}=1-\frac{610.69}{1495}=59.2\%.
\]
标准化后$T=(n-1)p=299\times5=1495$。59.2\%只表示：在这5个变量、这种标准化方式、$k=5$和这组初始点之下，类间平方和占总平方和的比例。换一组变量、换尺度或换$k$，这个数都会变；它本身并不说明分类“好”或“正确”，还要结合各类的分离程度和专业意义来评价（见下文）。
\subsection{类数的选择}\label{sec:k-number}
对固定的数据，最优类内平方和随$k$增大而不增：把某一类再拆成两类，$W$不会变大；$k=n$时$W=0$。因此$B/T$总是随$k$增大，不能单独用来确定类数，只能看曲线在哪里变平（“肘部”）。
\paragraph{轮廓系数} 对样品$i$，记$a(i)$为它到同类其他样品的平均距离，$b(i)$为它到其他各类样品平均距离中的最小值，
\[
s(i)=\frac{b(i)-a(i)}{\max\{a(i),b(i)\}}\in[-1,1].
\]
$s(i)$接近1表示分得清楚，接近0表示处在两类边界，为负表示可能分错。全部样品的平均轮廓系数可用来比较不同的$k$；经验上0.26～0.50表示结构较弱，0.51以上表示结构合理。
\paragraph{重抽样稳定性} (1) 对样品作有放回重抽样（或随机抽取80\%）；(2) 在重抽样数据上用同样的变量、标准化和$k$重新聚类；(3) 把每个原始类与重抽样结果中最相似的类配对，计算Jaccard系数（或调整兰德指数）；(4) 重复多次（如100次）取平均。平均Jaccard系数低于0.6的类不稳定，不宜作专业解释。R的\texttt{fpc::clusterboot()}可完成上述步骤。

耳测量例（标准化5变量，每个$k$用50套随机初始点）的复算结果如下：
\begin{center}\small\begin{tabular}{lccccccc}\toprule
{\bfseries $k$}&2&3&4&5&6&7&8\\\midrule
$W$&1039.7&835.4&694.2&610.7&545.0&495.9&457.4\\
$B/T$&30.5\%&44.1\%&53.6\%&59.2\%&63.5\%&66.8\%&69.4\%\\
平均轮廓系数&0.29&0.25&0.26&0.26&0.26&0.26&0.24\\\bottomrule
\end{tabular}\end{center}
$B/T$平稳上升，看不到明显的肘部；平均轮廓系数在0.24～0.29之间，说明耳朵形态更像一个连续分布，并没有界限分明的自然类别。按百分位数指定初始点的$k=5$解与50套随机初始点的最优解相同（$B/T=59.2\%$）。所以这里选$k=5$，主要依据是临床上需要几种代表性耳型，而不是数据中存在5个天然的类。
\begin{sikao}[$B/T$的参照：无结构数据的结果]
在单位正方形内均匀生成100个点，此数据不含任何群组结构，同样用K均值法分为3类；另生成一份确有3个群组的数据作为对照：
\par\noindent\begin{rcode}
library(cluster)
set.seed(3)
u <- matrix(runif(200),100,2)                     # 无结构：均匀分布
km <- kmeans(u,3,nstart=20)
km$betweenss/km$totss; mean(silhouette(km$cluster,dist(u))[,3])
v <- rbind(matrix(rnorm(70,0,.5),35), matrix(rnorm(70,3,.5),35),
           cbind(rnorm(30,0,.5),rnorm(30,3,.5)))  # 真有3个群
km2 <- kmeans(v,3,nstart=20)
km2$betweenss/km2$totss; mean(silhouette(km2$cluster,dist(v))[,3])
\end{rcode}
\begin{routput}
[1] 0.664
[1] 0.427
[1] 0.885
[1] 0.678
\end{routput}
无结构数据的$B/T$达66\%，高于耳测量资料分5类时的59.2\%。原因在于K均值法本身以最大化$B/T$为目标，总能找到使类间差异显得较大的划分。因此，$B/T$较高既不能说明分类效果好，也不能说明数据中存在真实的类别。

平均轮廓系数的区分能力较好：有真实结构时为0.68，无结构时为0.43。耳测量资料约为0.26，与前文“耳朵形态接近连续分布”的判断一致。该分类在临床上仍有使用价值（提供若干代表性耳型），但应视为对连续分布的人为划分，而非发现了天然存在的耳型。
\end{sikao}
\subsection{还原到原量纲：代表性耳朵}
标准化中心不便临床使用。把中心乘回各变量的标准差、加上均数，即还原到原量纲；再在每类中找离中心最近的实际病例，作为该类的“代表性耳朵”：
\begin{code}
km \ROperator{<-} \RFunction{kmeans}(db, ca)
mu \ROperator{<-} \RFunction{attr}(db, \RString{"scaled:center"}); sdv \ROperator{<-} \RFunction{attr}(db, \RString{"scaled:scale"})
\RFunction{sweep}(\RFunction{sweep}(km\ROperator{\$}centers, \RNumber{2}, sdv, \RString{"*"}), \RNumber{2}, mu, \RString{"+"})   \RComment{\# 原量纲中心}
rep.id \ROperator{<-} \RFunction{sapply}(\RNumber{1}\ROperator{:}\RNumber{5}, \RKeyword{function}(l) \{
  idx \ROperator{<-} \RFunction{which}(km\ROperator{\$}cluster \ROperator{==} l)
  idx[\RFunction{which.min}(\RFunction{rowSums}(\RFunction{sweep}(db[idx, , drop \ROperator{=} \RKeyword{FALSE}], \RNumber{2}, km\ROperator{\$}centers[l, ])\ROperator{\textasciicircum{}}\RNumber{2}))]
\})
da[rep.id, v]
\end{code}
\begin{center}\small\begin{tabular}{lcccccccccccc}\toprule
&&\multicolumn{5}{c}{\bfseries 原量纲类中心}&&\multicolumn{5}{c}{\bfseries 离中心最近的病例}\\\cmidrule(lr){3-7}\cmidrule(lr){9-13}
{\bfseries 类}&{\bfseries 例数}&EC&EK&EZ&EI&AI&{\bfseries 编号}&EC&EK&EZ&EI&AI\\\midrule
1&24&5.94&2.68&1.88&45.3&70.6&68&6.1&2.8&1.9&45.9&67.9\\
2&87&6.11&3.24&1.56&53.2&48.2&260&6.2&3.3&1.6&53.2&48.5\\
3&73&5.91&3.13&1.97&53.1&62.9&84&5.9&3.1&2.0&52.5&64.5\\
4&50&6.71&3.61&2.03&53.9&56.4&93&6.9&3.7&2.0&53.6&54.1\\
5&66&6.83&3.18&2.14&46.7&67.2&15&6.7&3.2&2.1&47.8&65.6\\\bottomrule
\end{tabular}\end{center}
EC、EK、EZ单位为cm，EI、AI单位为\%。各类大致为：1类短而窄、外展较大；2类中等大小、贴头（外展最小）；3类短、宽度中等、外展中等；4类长而宽；5类长而窄、外展较大。编号68、260、84、93、15这5例的实际耳朵即可作为修复时参照的代表。
\subsection{程序随机指定初始点}
直接输入$k$值。每次运行的结果可能不同。
\begin{code}
k \ROperator{<-} \RNumber{5}
g \ROperator{<-} \RFunction{kmeans}(db, k)
g\ROperator{\$}centers
k \ROperator{<-} \RNumber{4}
g \ROperator{<-} \RFunction{kmeans}(db, k)
g\ROperator{\$}centers
\end{code}
\par\addvspace{5pt}\begin{center}\begin{minipage}{\linewidth}\centering
\small
\captionof{table}{随机初始点得到的标化类别中心}
\begin{tabular}{llrrrrr}\toprule
\textbf{设置/运行}&\textbf{类别}&\textbf{EC}&\textbf{EK}&\textbf{EZ}&\textbf{EI}&\textbf{AI}\\\midrule
\multirow{5}{*}{$k=5$} & 1 & -0.47173490 & 0.38037882 & 0.57018830 & 0.86005170 & 0.28540470\\
 & 2 & -0.93000630 & -0.13683475 & -1.44698450 & 0.81911760 & -1.29576260\\
 & 3 & -0.66961030 & -1.33867575 & -0.04299810 & -0.76302790 & 0.74908390\\
 & 4 & 0.47183800 & 0.66283512 & -0.53548830 & 0.16830930 & -0.81491080\\
 & 5 & 1.06647380 & 0.09343839 & 0.93191100 & -0.90652000 & 0.78034750\\
\cmidrule(lr){1-7}
\multirow{4}{*}{$k=4$，第1次} & 1 & 1.23474240 & 1.00974607 & 0.23644540 & -0.20020600 & -0.30578870\\
 & 2 & 0.21512140 & -0.86994045 & 0.55697720 & -1.09478240 & 1.02540860\\
 & 3 & -0.55779840 & -0.07646648 & -1.03017810 & 0.47655890 & -0.92855330\\
 & 4 & -0.33233610 & 0.41458886 & 0.71851850 & 0.74744290 & 0.40214590\\
\cmidrule(lr){1-7}
\multirow{4}{*}{$k=4$，第2次} & 1 & -0.34752740 & 0.12352200 & -1.13844679 & 0.46927834 & -1.12862640\\
 & 2 & -0.97562040 & -0.96322620 & 0.06379143 & -0.04123529 & 0.60274110\\
 & 3 & 0.92162320 & -0.25996660 & 0.77040502 & -1.11960365 & 0.83796470\\
 & 4 & 0.43309520 & 1.06931640 & 0.48417739 & 0.62832488 & -0.12724190\\
\cmidrule(lr){1-7}
\multirow{4}{*}{$k=4$，第3次} & 1 & 0.73422560 & 1.18069880 & 0.61786590 & 0.44483660 & -0.06691626\\
 & 2 & -0.99670010 & -0.64423310 & 0.13653890 & 0.32425000 & 0.47130587\\
 & 3 & -0.30847770 & 0.12267940 & -1.10264660 & 0.42611670 & -1.09548297\\
 & 4 & 0.68654100 & -0.56818940 & 0.66733310 & -1.22008960 & 0.93892307\\
\bottomrule\end{tabular}
\end{minipage}\end{center}

\subsection{多套初始点}
\texttt{nstart=n}随机指定$n$套初始中心，每套分别迭代到收敛，取$W$最小的一套作为结果。
\Note{“迭代30遍”的说法不确切。\texttt{nstart = 30}是用30套初始中心各运行一次再择优；每次运行的迭代次数上限由\texttt{iter.max}控制（默认10）。}
\begin{code}
k \ROperator{<-} \RNumber{4}
g \ROperator{<-} \RFunction{kmeans}(db, k, nstart \ROperator{=} \RNumber{30})
g\ROperator{\$}centers
\end{code}
\par\addvspace{5pt}\begin{center}\begin{minipage}{\linewidth}\centering
\small
\captionof{table}{\texttt{nstart=30} 的三次运行结果}\label{tab:nstart}
\begin{tabular}{llrrrrr}\toprule
\textbf{设置/运行}&\textbf{类别}&\textbf{EC}&\textbf{EK}&\textbf{EZ}&\textbf{EI}&\textbf{AI}\\\midrule
\multirow{4}{*}{第1次} & 1 & 0.92162320 & -0.25996660 & 0.77040502 & -1.11960365 & 0.83796470\\
 & 2 & -0.34752740 & 0.12352200 & -1.13844679 & 0.46927834 & -1.12862640\\
 & 3 & 0.43309520 & 1.06931640 & 0.48417739 & 0.62832488 & -0.12724190\\
 & 4 & -0.97562040 & -0.96322620 & 0.06379143 & -0.04123529 & 0.60274110\\
\cmidrule(lr){1-7}
\multirow{4}{*}{第2次} & 1 & 0.92162320 & -0.25996660 & 0.77040502 & -1.11960365 & 0.83796470\\
 & 2 & -0.97562040 & -0.96322620 & 0.06379143 & -0.04123529 & 0.60274110\\
 & 3 & 0.43309520 & 1.06931640 & 0.48417739 & 0.62832488 & -0.12724190\\
 & 4 & -0.34752740 & 0.12352200 & -1.13844679 & 0.46927834 & -1.12862640\\
\cmidrule(lr){1-7}
\multirow{4}{*}{第3次} & 1 & 0.43309520 & 1.06931640 & 0.48417739 & 0.62832488 & -0.12724190\\
 & 2 & -0.34752740 & 0.12352200 & -1.13844679 & 0.46927834 & -1.12862640\\
 & 3 & 0.92162320 & -0.25996660 & 0.77040502 & -1.11960365 & 0.83796470\\
 & 4 & -0.97562040 & -0.96322620 & 0.06379143 & -0.04123529 & 0.60274110\\
\bottomrule\end{tabular}
\end{minipage}\end{center}
表\ref{tab:nstart}三次运行得到的是同一组4个中心，只是类别编号不同。例如中心$(0.92,\allowbreak-0.26,\allowbreak0.77,\allowbreak-1.12,\allowbreak0.84)$在第1、2次编为1类，在第3次编为3类。类别编号本身没有意义，对编号作置换不改变分组，$W$与$B/T$（53.6\%）也完全相同。比较不同运行的结果时，应按中心或成员配对，而不是按编号。相比之下，不设\texttt{nstart}的三次$k=4$运行中，只有第2次达到这一最优解，第1、3次停在了不同的局部最优。

\begin{chaptersummary}{本章小结：聚类分析}
研究目的&在没有类别标签的情况下，把相似的样品（Q型）或指标（R型）归类，属于探索性分析\\
优化准则&层次聚类：按选定的类间距离逐步合并（Ward法每步使$W$增加最少）；K均值：最小化类内平方和$W$，等价于最大化$B/T$\\
主要条件&相似性度量要与变量类型匹配；数值变量量纲不同须先标准化；马氏距离要求协方差阵可逆；K均值适合大致球形、大小相近的类，须预先给定$k$\\
结果解释&聚类树、类中心（应还原到原量纲）、各类代表个体；结合$B/T$曲线、轮廓系数、重抽样稳定性与专业意义判断类数和效果\\
常见误用&把相似度当距离（配合例0.6应为距离0.4）；不分对称与不对称二值变量；以为\texttt{dist()}会自动标准化；只凭$B/T$定类数；按编号比较不同运行的结果；把任何数据都“分出来”的类当作天然存在的类别\\
\end{chaptersummary}

\begin{fuxi}[复习题]
\begin{enumerate}
\item 两个变量的单位分别为cm和kg，说明直接计算欧氏距离存在的问题，以及标准化所隐含的假定。\\
\textit{答：方差大的变量主导距离；标准化后各变量对平方距离的平均贡献相同，即假定各变量同等重要。}
\item 比较马氏距离与标准化欧氏距离。\\
\textit{答：标准化欧氏距离仅消除量纲（方差）的影响；马氏距离同时考虑变量间的相关，相当于先作白化变换再计算欧氏距离。}
\item 两名患者均无某种罕见症状，比较简单配合系数与Jaccard系数对这一情况的处理。\\
\textit{答：简单配合系数将(0,0)计为相同；Jaccard系数排除(0,0)，适用于仅阳性取值有意义的不对称二值变量。}
\item 说明最短距离法容易出现的现象，以及Ward法所得类别的特点。\\
\textit{答：最短距离法易出现链式效应；Ward法倾向于得到大小较均衡、近似球形的类。}
\item 说明K均值聚类需尝试多套初始点的原因，以及$B/T$在确定类数中的作用。\\
\textit{答：K均值法只能保证收敛到局部最优，结果依赖初始点。$B/T$随$k$增大而不减，只能观察其“肘部”，类数应结合轮廓系数、稳定性与专业需要确定。}
\end{enumerate}
\end{fuxi}
